\chapter{PWM 控制技术}

PWM（Pulse Width Modulation，脉冲宽度调制）技术是电力电子技术中最重要的控制技术之一。它没有引入新的电路拓扑，却使电路的控制性能大为改善，凡是用全控型器件的电路（逆变、整流、斩波、交交变频）都可以用它。本章讲四件事：PWM 的理论依据、PWM 逆变电路及其控制方法、跟踪型 PWM、PWM 整流电路。

% ============================================================
\section{PWM 控制的基本原理}

PWM 控制的理论基础是\textbf{面积等效原理}（冲量等效原理）：

\begin{tcolorbox}[enhanced,breakable,colback=highlight,colframe=yhbAccent!30!white,boxrule=0.5pt,arc=2.5pt,left=11pt,right=11pt,top=9pt,bottom=9pt,boxsep=0pt,before skip=11pt,after skip=11pt]
\textbf{面积等效原理}\quad 冲量相等而形状不同的窄脉冲加在具有惯性的环节上时，其效果基本相同。“冲量”指窄脉冲的面积，“效果基本相同”指环节的输出响应波形基本相同。
\end{tcolorbox}

由此可以推得：把正弦半波分成 $N$ 等份，用 $N$ 个\textbf{等幅不等宽}的矩形脉冲代替，使每个矩形脉冲的\textbf{面积与对应正弦波段的面积相等}，则这串脉冲加到惯性环节上，其效果与正弦半波基本相同。这就是 \textbf{SPWM}（正弦波脉宽调制）的由来。

\begin{itemize}
  \item 脉冲的\textbf{宽度}按正弦规律变化、\textbf{幅值}恒定——称为 PWM 波形；
  \item 同样的道理，也可以把非正弦的期望波形（如梯形波、指定谐波消除波形）用等幅不等宽脉冲等效；
  \item 用 PWM 波形代替正弦波，关键在于\textbf{开关频率足够高}（远高于输出基波频率）时，负载（电动机等惯性环节）只“感受到”平均效果。
\end{itemize}

\corepoint{PWM 控制的三个基本量必须分清：\textbf{调制波}（期望输出的参考波，通常是正弦波）、\textbf{载波}（高频三角波或锯齿波，决定开关频率）、\textbf{调制比} $m=U_{rm}/U_{cm}$（决定输出基波幅值）、\textbf{载波比} $N=f_c/f_r$（决定每周期内的脉冲个数与谐波分布）。PWM 的输出电压基波幅值与 $m$ 成正比，这是它“调压”的原理；而开关频率决定了谐波往高频搬移，这是它“改善波形”的原理。}

% ============================================================
\section{PWM 逆变电路及其控制方法}

\subsection{计算法和调制法}

\textbf{计算法}：根据预期的输出电压波形，按面积等效原理计算出各脉冲的宽度与间隔，用数字方式（查表或实时计算）生成 PWM 波形。计算法复杂、不灵活，但可以\textbf{精确控制}波形（如指定谐波消除），多用于离线计算、查表运行的场合。

\textbf{调制法}：把期望的输出波形作为\textbf{调制信号}，把接受调制的信号作为\textbf{载波}，通过比较得到 PWM 波形。这是工程上应用最广的方法。

\begin{itemize}
  \item \textbf{单极性调制}：载波为单极性三角波（只在正或负半周变化），调制波为正弦，比较结果为\textbf{单极性}脉冲列（$0$ 或 $+U_d$，或 $0$ 或 $-U_d$），输出谐波少、但需两只管子按半周交替工作。
  \item \textbf{双极性调制}：载波为双极性三角波（在正负之间变化），比较结果为\textbf{双极性}脉冲列（$+U_d$ 或 $-U_d$），控制简单、同一桥臂上下管互补，是三相桥 SPWM 最常用的方式。
\end{itemize}

单相全桥双极性 SPWM 的输出电压基波幅值为

\keyformula{U_{o1m}=mU_d\qquad(0\le m\le1)}

三相桥式 SPWM 逆变电路（三个桥臂共用同一三角载波，三相调制波互差 $120^\circ$）输出相电压的基波幅值为

\keyformula{U_{UN1m}=\frac{mU_d}{2}\qquad(0\le m\le1)}

\subsection{异步调制和同步调制}

按载波比 $N=f_c/f_r$ 是否固定，分为两种调制方式：

\begin{table}[htbp]
\centering
\caption{异步调制与同步调制}
\label{tab:7-1}
\begin{tabular}{p{2.4cm}p{5.4cm}p{6.2cm}}
\toprule
\textbf{方式} & \textbf{做法} & \textbf{特点} \\
\midrule
异步调制 & 载波频率 $f_c$ 固定，$N$ 随 $f_r$ 变化 & 低频时 $N$ 很大、脉冲多、波形好；高频时 $N$ 变小、正负半周脉冲不对称，产生偶次谐波与波形畸变 \\
同步调制 & $N$ 固定（$f_c$ 随 $f_r$ 成正比变化） & 正负半周严格对称、无偶次谐波，但低频时 $N$ 太小、谐波次数低 \\
分段同步调制 & 把 $f_r$ 分成若干段，每段内取不同的固定 $N$ & 兼顾两者：低频段用大 $N$、高频段用小 $N$，是实际变频器常用方案 \\
\bottomrule
\end{tabular}
\end{table}

\subsection{规则采样法}

自然采样法的脉宽由正弦波与三角波的自然交点决定，需解超越方程，难以实时计算。\textbf{规则采样法}用固定的采样点近似：

\begin{itemize}
  \item \textbf{对称规则采样}：在每个载波周期内的固定时刻（如三角波顶点）对调制波采样，以该值决定脉宽，脉宽关于载波周期中点对称；
  \item \textbf{不对称规则采样}：在三角波的顶点和底点各采样一次，脉宽不对称，精度更高。
\end{itemize}

对称规则采样法的脉宽公式为

\keyformula{\delta=\frac{T_c}{2}\left(1+m\sin\omega_r t_s\right)}

式中 $T_c$ 为载波周期，$t_s$ 为采样时刻。规则采样法计算简单、便于数字实现，其输出与自然采样法的差别在开关频率足够高时很小。

\subsection{PWM 逆变电路的谐波分析}

PWM 逆变电路输出电压的谐波分布与载波比 $N$ 密切相关。以三相桥双极性 SPWM 为例（$m$ 一定、$N$ 为整数）：

\begin{itemize}
  \item 输出线电压中\textbf{不含低次谐波}，最低次谐波出现在\textbf{载波频率 $f_c$ 附近}（即 $N$ 次及其边频 $N\pm2$、$N\pm4$ 等）；
  \item 谐波幅值随 $m$ 增大而增大，但\textbf{谐波次数不随 $m$ 改变}；
  \item $N$ 越大，最低次谐波次数越高、越容易被滤除，故提高载波频率是改善波形的主要手段（代价是开关损耗增大）。
\end{itemize}

\subsection{提高直流电压利用率和减少开关次数}

SPWM 的\textbf{直流电压利用率低}是它的主要缺点：$m\le1$ 时三相桥输出线电压基波幅值最大为 $\dfrac{\sqrt{3}}{2}U_d=0.866U_d$，即\textbf{利用率只有 0.866}。提高的办法有：

\begin{enumerate}
  \item \textbf{注入三次谐波}：在正弦调制波中注入幅值为基波 $1/6$ 的三次谐波（三线制中三次谐波不产生线电流），可使调制波峰值降低而不改变基波，从而把直流电压利用率\textbf{提高约 15\%}；
  \item \textbf{梯形波调制}：用梯形波作调制波，利用率可提高到 1 左右，但输出含较多低次谐波；
  \item \textbf{空间矢量 PWM（SVPWM）}：从电动机磁链圆的角度出发选择电压矢量，直流电压利用率可达 \textbf{1.0}（比 SPWM 高约 15.5\%），且便于数字实现，是目前变频器的主流算法。
\end{enumerate}

\keyformula{\text{SPWM: }\frac{U_{UV1m}}{U_d}\le0.866;\qquad \text{SVPWM: }\frac{U_{UV1m}}{U_d}\le1.0}

\subsection{空间矢量 PWM 控制}

三相桥式逆变电路的 6 个开关共有 8 种合法开关状态，对应 \textbf{8 个基本电压矢量}：6 个幅值为 $\frac{2}{3}U_d$ 的非零矢量（$U_1\sim U_6$，互差 $60^\circ$，构成正六边形）与 2 个零矢量（$U_0$、$U_7$）。

\begin{itemize}
  \item 期望的电压矢量落在某个 $60^\circ$ 扇区内时，可用该扇区相邻的两个非零矢量与零矢量\textbf{在一个载波周期内分时合成}（矢量合成即“伏秒平衡”）；
  \item 设两相邻矢量作用时间为 $T_1$、$T_2$，零矢量作用时间为 $T_0$，则 $T_1+T_2+T_0=T_c$，由伏秒平衡可解出 $T_1$、$T_2$；
  \item 优点：直流电压利用率高、开关次数少（每次只切换一个桥臂）、便于微机实时计算。
\end{itemize}

\subsection{PWM 逆变电路的多重化}

把几组 PWM 逆变电路按一定相位差组合（多重化），可使载波频率附近的谐波相互抵消、等效开关频率成倍提高，从而在\textbf{不提高器件开关频率}的前提下改善波形——这与第 4 章的多重化思想一致，只是这里抵消的是 PWM 的高频谐波。

% ============================================================
\section{PWM 跟踪控制技术}

\textbf{跟踪控制}不预先计算脉冲宽度，而是把期望输出（电流或电压）作为指令，\textbf{实时比较实际输出与指令的偏差}，由偏差直接决定开关动作，使实际输出跟踪指令变化。它属于“闭环 PWM”。

\subsection{滞环比较方式}

把指令信号与实际信号送入\textbf{滞环比较器}（环宽为 $\pm\Delta I$），当偏差超出环宽时改变开关状态。以电流跟踪型 PWM 为例：

\begin{itemize}
  \item 实际电流上升超过上限 $I^*+\Delta I$ 时，关断使电流下降的开关组合；低于下限 $I^*-\Delta I$ 时，开通使电流上升的组合；
  \item \textbf{环宽越小}，跟踪误差越小，但开关频率越高、开关损耗越大；\textbf{环宽越大}，开关频率低而误差大；
  \item 优点：硬件简单、响应快、不需要载波，天然限流；缺点：\textbf{开关频率不固定}（随电流变化率而变），给滤波设计带来困难，且输出谐波频谱较宽。
\end{itemize}

\subsection{三角波比较方式}

把指令电流与实际电流的偏差经 PI 调节后，与\textbf{三角波载波}比较产生 PWM 信号。

\begin{itemize}
  \item \textbf{开关频率固定}（等于载波频率），便于滤波器设计，输出谐波集中在载波频率附近；
  \item 由于有 PI 调节器，稳态误差小、跟踪精度高；
  \item 缺点是响应速度略慢于滞环方式，且 PI 参数需整定。
\end{itemize}

% ============================================================
\section{PWM 整流电路及其控制方法}

\subsection{PWM 整流电路的工作原理}

PWM 整流电路与 PWM 逆变电路结构相同（电压型／电流型、单相／三相），只是\textbf{能量流向相反}：它把交流电变成直流电，但通过 PWM 控制可使\textbf{交流侧电流接近正弦}且\textbf{功率因数可调（可做到单位功率因数，甚至电流超前）}，还能实现\textbf{能量的双向流动}（可整流也可逆变）——这是相控整流做不到的。

\begin{itemize}
  \item \textbf{电压型 PWM 整流电路}：直流侧为大电容，交流侧经电感接入电网，桥臂由全控器件与反并联二极管构成；
  \item \textbf{电流型 PWM 整流电路}：直流侧为大电感，交流侧需电容滤波。
\end{itemize}

工作原理要点：通过控制桥臂开关，使桥臂中点电压（PWM 电压）中含有与电网电压同频同相的基波分量，控制该基波的\textbf{幅值与相位}，即可控制交流侧电流的大小与相位——本质上是用 PWM 电压“顶”着电网电压，靠电感上的电压差控制电流。

\subsection{PWM 整流电路的控制方法}

\begin{enumerate}
  \item \textbf{间接电流控制}（幅相控制）：直接控制 PWM 电压的幅值与相位，间接控制交流电流。结构简单、无需电流传感器，但\textbf{动态响应慢}、对参数敏感，电流波形受电网电压畸变影响大。
  \item \textbf{直接电流控制}：引入交流电流闭环，把电流指令与实际电流的偏差经调节器后生成 PWM 指令。又分为：
    \begin{itemize}
      \item \textbf{滞环电流控制}：环宽决定开关频率，响应快；
      \item \textbf{固定开关频率的电流控制}（如三角波比较、同步旋转坐标系 PI 控制）：开关频率固定、电流正弦度好，工业上最常用；
      \item \textbf{预测电流控制}：按下一拍期望电流反算所需电压，动态最快。
    \end{itemize}
  \item 工程上电压型 PWM 整流通常采用\textbf{电压外环 + 电流内环}的双闭环结构：电压外环稳定直流母线电压并给出有功电流指令，电流内环实现快速跟踪；无功电流指令置零即得单位功率因数。
\end{enumerate}

\corepoint{本章必须记牢的数字与概念：\textbf{调制比 $m=U_{rm}/U_{cm}$、载波比 $N=f_c/f_r$}；\textbf{单相全桥 $U_{o1m}=mU_d$、三相桥相电压 $U_{UN1m}=mU_d/2$}；\textbf{SPWM 直流电压利用率 0.866、SVPWM 为 1.0（提高 15.5\%）、三次谐波注入提高 15\%}；\textbf{异步调制低频好高频差、同步调制对称但低频谐波大、分段同步折中}；\textbf{滞环比较开关频率不固定、三角波比较频率固定}；\textbf{PWM 整流能实现单位功率因数与能量双向流动}。}
